% Uses the macros and colours of the preamble (\nccl, \bine, \pat, \todo, ch0, ch1, ...).
\chapter{PAT-Based Implementation}\label{ch:pat}

This chapter describes a second implementation of the \bine{} allgather in \nccl{}.
It does not change the schedule of \cref{ch:butterfly}: the partners, the labels, the four strategies for the non-contiguous data and the invariants are the same.
What changes is how the schedule is executed.
The kernels of \cref{ch:butterfly} issue one primitive call per transfer, from two groups of threads; here the allgather uses the device-side machinery of \pat{} (\cref{sec:nccl-pat}), a schedule thread that feeds a set of workers, which pack many pieces in each slot of the FIFO.
Only the allgather is implemented, for the \texttt{Simple} protocol, the only one for which \pat{} exists in \nccl{}.

%==============================================================================
\section{Motivation}\label{sec:pat-motivation}

The kernels of \cref{sec:bfly-kernels} send and receive from separate thread groups and fuse the copy of the own block into the first send, so what remains of their cost depends on the strategy and on the layout of the data, not on the schedule.
\begin{itemize}
  \item \emph{Number of calls.} Every call consumes the FIFO steps of one chunk, whatever its size (\cref{sec:bfly-proxy}).
  A window that is not one contiguous run is transferred with one call per block and per chunk (strided transfers, \cref{sec:bfly-kernels}), so a small piece occupies a whole FIFO step. This is always the case for block-by-block, and for every strategy in slice mode when the slice does not cover whole blocks.
  \item \emph{Partition.} Pos mode (\cref{sec:bfly-pos}) removes the stride for small messages, but not for permute, and not for messages above the threshold $S$, where strided transfers return.
  \item \emph{Per-step overhead.} Two \texttt{Primitives} objects are built at every step, and the steps are separated by a barrier of the whole block.
\end{itemize}
The aggregation of \pat{} addresses the first two.
A slot of the FIFO is filled by many operations, each reading a piece from an arbitrary place of the output buffer, and is posted once.
The non-contiguity of the blocks is absorbed by the local copy into the slot, executed by the workers, and the network sees one message per slot whatever the strategy and the partition.
The price is one more local copy for every received piece (\cref{sec:pat-prims}).

%==============================================================================
\section{Mapping the schedule on \pat{}}\label{sec:pat-mapping}

\paragraph{Same tables.}
The kernels read the tables of \cref{ch:butterfly} from \texttt{channel.bine}: $\Pi$ for send, permute and block-by-block, $\Pi^{\mathrm{dh}}$ for double send, and the labels $\nu$ and $\nu^{-1}$.
Step $k\in[0,s-1]$ is the same as in \cref{alg:bfly-allgather}: the partner is $\Pi(r,s-1-k)$, and for double send $\Pi^{\mathrm{dh}}(r,s-1-k)=\Pi(r,k)$ (\cref{sec:bfly-dhlv}).
The window has $2^k$ blocks and doubles at each step; the slots of the output buffer are those of \cref{tab:bfly-strategies}.

\paragraph{Steps as dimensions.}
The dimension $d$ of \pat{} is the distance $2^d$ between peers, the same for every rank.
In the \bine{} butterfly the offset depends on the rank, but each step still identifies one peer per rank, and the \pat{} structures only need an index for it.
We therefore use the step as the dimension: the connector of dimension $k$ of a rank points to its partner of step $k$.
The peer to receive from and the peer to send to are the same rank, since the butterfly is symmetric; this is the only change in the constructor of the primitives (\cref{sec:pat-prims}).

\paragraph{No one-rank-per-node restriction.}
The native \pat{} is limited to one rank per node (\cref{sec:nccl-pat}); this restriction is not inherited.
The peers are given by the tables of \cref{ch:butterfly}, which depend only on rank numbers, and the \bine{} setup selects for each pair the transport that can connect it (\cref{sec:nccl-transport}), so a partner can be on the same node, on another node, or both.

\paragraph{Partition.}
The implementation uses slice mode only: a channel handles a slice of every block, and pos mode (\cref{sec:bfly-pos}) is not used.
For small messages the aggregation plays its role, since the pieces of many blocks share a slot.

\paragraph{Comparison with \pat{}.}
\Cref{tab:pat-vs-bine} compares the two.
\bine{} replaces the Bruck-like binomial trees with a butterfly, so each step exchanges the whole window with one partner and there is no linear phase of its own: when the window is larger than a slot, the step is split in more slots, which plays the role of the linear part.
As in \pat{}, the steps with the most data are those with the closest partners, and \bine{} reduces these distances further (\cref{sec:bg-bine-butterfly}).

\begin{table}[t]
\centering
\footnotesize
\setlength{\tabcolsep}{4pt}
\caption{Original \pat{} allgather and the \bine{} allgather of this chapter.}
\label{tab:pat-vs-bine}
\begin{tabularx}{\linewidth}{@{}l X X@{}}
\toprule
 & \pat{} (\nccl{}) & \bine{} on \pat{} machinery \\
\midrule
Pattern & aggregated binomial trees, any $p$ & \bine{} butterfly, $p=2^s$ \\
Peer of dimension $d$ & rank at distance $2^d$ & partner of step $d$ (\cref{ch:butterfly}) \\
Local data & kept in the slot and forwarded without copy & always in the output buffer, sent with \texttt{PatSrcOutput} \\
Operation & receive and send can be fused in one step & each operation either sends or receives \\
Aggregation & limited by the buffer, then linear phase & a step spans $N_k$ slots (\cref{eq:pat-slots}) \\
Non-contiguous data & not applicable & four strategies (\cref{tab:pat-strategies}) \\
Partition & slice of every block & slice mode only \\
GPUs per node & one (inter-node part only) & any \\
Protocols, zero-copy & \texttt{Simple}, no zero-copy & same, inherited from \pat{} \\
\bottomrule
\end{tabularx}
\end{table}

%==============================================================================
\section{Geometry: pieces and slots}\label{sec:pat-geom}

A channel handles the slice $[\texttt{channelOffset},\texttt{channelOffset}+\texttt{channelCount})$ of every block (\cref{sec:nccl-cbd}).
The slice is cut in chunks, and every operation moves a \emph{piece}, the part of one chunk of one block.
Let $E_s$ be the number of elements of a slot (\texttt{slotElems}), i.e., the size of the \texttt{Simple} buffer divided by \texttt{NCCL\_STEPS} (\cref{sec:nccl-transport}).
For the allgather an element is a byte, and in the default configuration the buffer is $4$\,MiB and a slot $512$\,KiB~\cite{hu2025demystifying}.
The structure \texttt{BineAgGeom} derives, from the values of \texttt{ncclCollCbdPart},
\begin{align}
  c &= \min(\texttt{chunkCount},\,E_s), & n_c &= \lceil \texttt{channelCount}/c\rceil,\\
  \sigma &= \min(\texttt{channelCount},\,c), & P &= \max\bigl(1,\lfloor E_s/\sigma\rfloor\bigr),
\end{align}
where $c$ is the chunk (never larger than a slot), $n_c$ the number of chunks, $\sigma$ the \emph{stride}, i.e., the space reserved to a piece in the slot, and $P$ the number of pieces per slot.

At step $k$ the window has $2^k$ blocks, so the step moves $n_c\,2^k$ pieces: piece $q\in[0,n_c2^k-1]$ is chunk $\lfloor q/2^k\rfloor$ of block $j=q\bmod 2^k$ of the window.
The pieces are grouped $P$ at a time, so the step uses
\begin{equation}
  N_k=\Bigl\lceil \frac{n_c\,2^k}{P}\Bigr\rceil
  \label{eq:pat-slots}
\end{equation}
slots, all full except possibly the last, and piece $q$ is placed at offset $(q\bmod P)\,\sigma$ of its slot (\cref{fig:pat-slots}).
Sender and receiver compute the same layout, since $n_c$, $P$ and $k$ depend only on the work descriptor and the step, and the order of the pieces is the order of the blocks of the window, on which the two partners agree (invariant~\ref{inv:bfly-split}).
The pre-step of \emph{send} uses one slot per chunk.

\begin{figure}[t]
\centering
\begin{tikzpicture}[x=1.15cm, y=-0.95cm, font=\footnotesize,
    cell/.style={draw, minimum width=1.05cm, minimum height=0.7cm, inner sep=0pt}]
  \foreach \sl/\u/\lab/\col in {0/0/0/ch0, 0/1/1/ch0, 0/2/2/ch0,
                               1/0/3/ch0, 1/1/0/ch1, 1/2/1/ch1,
                               2/0/2/ch1, 2/1/3/ch1}
    \node[cell, fill=\col!35] at (\u,\sl) {$j{=}\lab$};
  \node[cell, dashed, draw=black!50] at (2,2) {};
  \foreach \sl in {0,1,2} {
    \draw[rounded corners=2pt] (-0.62,\sl-0.5) rectangle (2.62,\sl+0.5);
    \node[anchor=east] at (-0.8,\sl) {slot $g_i{=}\sl$};
    \node[anchor=west] at (2.85,\sl) {one message};
  }
  % legend
  \node[cell, fill=ch0!35, minimum width=0.5cm, minimum height=0.4cm] at (0,3.1) {};
  \node[anchor=west] at (0.3,3.1) {chunk 0};
  \node[cell, fill=ch1!35, minimum width=0.5cm, minimum height=0.4cm] at (1.8,3.1) {};
  \node[anchor=west] at (2.1,3.1) {chunk 1};
\end{tikzpicture}
\caption{Packing of the pieces of a step into slots, with a window of $2^k=4$ blocks, $n_c=2$ chunks and $P=3$ pieces per slot: $N_k=\lceil 8/3\rceil=3$ slots, the last one not full. Pieces are ordered chunk by chunk; $j$ is the block of the window.}
\label{fig:pat-slots}
\end{figure}

Two regimes follow from \cref{eq:pat-slots}.
For small messages the slice is a single chunk, $\sigma=\texttt{channelCount}$ and $P=\lfloor E_s/\texttt{channelCount}\rfloor$ is large: as long as $2^k\le P$ the step is one slot, hence one transfer, and the algorithm is latency-optimal with $s$ transfers.
For large messages $P=1$ and a step moves $n_c2^k$ slots; this is the linear part, which moves full slots and, per channel, about $(p-1)\,\texttt{channelCount}/E_s$ slots over all the steps.

%==============================================================================
\section{Schedule generator}\label{sec:pat-gen}

\subsection{Operation order}

The schedule thread asks a generator, a class with the method \texttt{getNextOp}, for one operation at a time and writes it in an entry of a ring (\cref{sec:nccl-pat-impl}).
There is one class per strategy (\texttt{BineAGAlgorithm} for send, \texttt{BineAGPermutationAlgorithm}, \texttt{BineAGBlockByBlockAlgorithm}, \texttt{BineAGDoubleSendAlgorithm}), sharing the structure of \cref{alg:pat-gen}.
The generator is a state machine whose state is the stage (pre-step, butterfly, end), the step $k$, the slot $g_i$ and the position inside the slot, so that it produces one operation per call.

The order of the operations is the order in which each connection is used, and it must be the one the proxy expects (\cref{sec:nccl-proxy}); the proxy is the one of \cref{ch:butterfly}, unchanged.
In the pre-step, for every chunk, a send and then a receive; in a step, for every slot, the send of $P$ pieces (posted at the last one), then the receive of $P$ pieces (posted at the last one).
Unlike the kernels of \cref{sec:bfly-kernels}, where sends and receives are executed by different thread groups and progress independently, one group of workers executes all the operations of a channel in order.
Sends and receives must therefore alternate slot by slot, so that no rank blocks in a send that its partner cannot yet drain.

\begin{algorithm}[t]
\caption{Operations emitted by the schedule thread (common structure of the four strategies).}
\label{alg:pat-gen}
\begin{algorithmic}[1]
\ForAll{chunks $e$ of the slice}\Comment{pre-step, strategy dependent}
  \State emit \Call{PreOp}{$e$}
\EndFor
\For{$k\gets0$ \textbf{to} $s-1$}
  \State $(\mathit{sendWin},\mathit{recvWin})\gets$ \Call{Windows}{$k$} \Comment{\cref{tab:pat-strategies}}
  \For{$g_i\gets0$ \textbf{to} $N_k-1$}
    \State $\ell\gets\min(P,\ n_c2^k-g_iP)$
    \For{$u\gets0$ \textbf{to} $\ell-1$}\Comment{send the slot}
      \State $q\gets g_iP+u$;\quad emit \texttt{send}$(\textsc{Block}(\mathit{sendWin},q\bmod2^k),\ \textsc{Chunk}(q),\ \text{offset}=u\sigma)$
      \State \textbf{if} $u=\ell-1$ \textbf{then} mark the operation as \emph{postSend}
    \EndFor
    \For{$u\gets0$ \textbf{to} $\ell-1$}\Comment{receive the slot}
      \State $q\gets g_iP+u$;\quad emit \texttt{recv}$(\textsc{Block}(\mathit{recvWin},q\bmod2^k),\ \textsc{Chunk}(q),\ \text{offset}=u\sigma)$
      \State \textbf{if} $u=\ell-1$ \textbf{then} mark the operation as \emph{postRecv}
    \EndFor
  \EndFor
\EndFor
\State emit the end marker \Comment{$\texttt{last}=2$}
\end{algorithmic}
\end{algorithm}

Each operation has \texttt{sendDim} or \texttt{recvDim} equal to the dimension of the step, the offset \texttt{channelOffset} plus that of the chunk, and \texttt{nelem} equal to the length of the piece (shorter for the last chunk).
A channel with no data emits only the end marker, with the \texttt{PatSkipped} flag.

\subsection{Strategies}\label{sec:pat-strategies}

The four generators differ only in the choices of \cref{tab:bfly-strategies}: pre-step, slot of a piece, post-step.
\Cref{tab:pat-strategies} gives them in terms of the generator, with $x=\mathrm{pos}(r)$ (\texttt{myIdx} in the code) and $\textsc{Block}(w,j)$ the block of the output buffer that piece $j$ of a window starting at $w$ reads from or writes to.

\begin{table}[t]
\centering
\footnotesize
\setlength{\tabcolsep}{4pt}
\caption{Choices of the generator for the four strategies. $g=2^k$, $b=\lfloor x/2g\rfloor 2g$.}
\label{tab:pat-strategies}
\begin{tabularx}{\linewidth}{@{}l X X X X@{}}
\toprule
Strategy & Pre-step (per chunk) & Windows at step $k$ & $\textsc{Block}(w,j)$ & Post-step \\
\midrule
Send & exchange: send my block, receive into slot $x$; or local copy to slot $r$ if the peer is $r$ & $b$ or $b+g$ according to bit $k$ of $x$ & $w+j$ (position) & -- \\
Permute & local copy to slot $x$ & as send & $w+j$ & in-place permutation, block $b\leftarrow$ slot $\mathrm{pos}(b)$ \\
Block-by-block & local copy to slot $r$ & as send & $\nu^{-1}(\mathrm{rev}(w+j))$ & -- \\
Double send & local copy to slot $r$ & circular window of $g$ blocks, adjacent received half & $(w+j)\bmod p$ & -- \\
\bottomrule
\end{tabularx}
\end{table}

\begin{itemize}
  \item \emph{Send.} The window is the aligned range of position $x$; the half sent is the one the rank holds (the lower one if bit $k$ of $x$ is $0$, the upper one otherwise) and the other is received. Blocks are read and written at the slot of their position.
  The pre-step moves the block of each rank to its position: if the rank is its own destination (\texttt{redistTo}$=r$) it is a local copy; otherwise the rank sends its block to $\nu^{-1}(\mathrm{rev}(r))$ and receives into slot $x$ the block of rank $x$. The send is flagged \texttt{PatNoLocalCopy}, since it reads the input buffer and must not write slot $r$.
  \item \emph{Permute.} Blocks are placed at their position with a local copy; the butterfly is the same as send, and the permutation to natural order runs after the schedule (\cref{sec:pat-permute}).
  \item \emph{Block-by-block.} Blocks stay in natural order, so piece $j$ of a window is block $\nu^{-1}(\mathrm{rev}(w+j))$; the windows are the same as send. This is the strategy that benefits most from aggregation: the blocks are not adjacent, but the workers pack the pieces in the slot and they travel in one message.
  \item \emph{Double send.} The generator keeps a circular window starting as $\{r\}$ (\texttt{winStart}$=r$). At step $k$ it sends the current window and receives the adjacent half of size $2^k$, below it if $r$ and $k$ have different parity and above it otherwise, the time reverse of the rule of \cref{sec:bfly-dhlv}; the received half is merged in the window. All indices are taken modulo $p$ with the mask $p-1$. The dimension of step $k$ is $s-1-k$ of $\Pi^{\mathrm{dh}}$, i.e., the partner $\Pi(r,k)$, and no pre-step or post-step is needed.
\end{itemize}

%==============================================================================
\section{Kernel}\label{sec:pat-kernel}

The kernel \texttt{runBineAgPat} is that of \pat{}, instantiated with the generator of the strategy.
\Cref{alg:pat-kernel} shows it, and \cref{fig:pat-roles} the roles of the threads.
The geometry is derived once from the descriptor, the shared ring is cleared, and the threads are split in three roles: the schedule thread, the workers, and the remaining threads, which only take part in the barriers.
All the threads reach the same point at the end, where the strategy can run a hook that uses \texttt{\_\_syncthreads} (the post-step of permute).

\begin{figure}[t]
\centering
\begin{tikzpicture}[node distance=1.7cm, >=Stealth, font=\footnotesize,
    box/.style={draw, rounded corners=2pt, minimum height=1cm, minimum width=2.2cm, align=center, fill=white}]
  \node[box, fill=ch0!15] (gen)  {schedule thread\\ \texttt{getNextOp}};
  \node[box, right=of gen] (ring) {ring\\ \texttt{patSteps}};
  \node[box, fill=ch1!15, right=of ring] (wrk) {workers\\ \texttt{patCopy}};
  \node[box, right=of wrk] (net) {FIFO slots\\ dimension $d$};
  \draw[->, thick] (gen) -- node[above=1pt] {publish} (ring);
  \draw[->, thick] (ring) -- node[above=1pt] {poll} (wrk);
  \draw[->, thick] (wrk) -- node[above=1pt] {send/recv} (net);
  \node[below=0.5cm of ring, font=\footnotesize] {remaining threads: barriers only};
\end{tikzpicture}
\caption{Roles of the threads of a channel in \texttt{runBineAgPat}. The schedule thread runs ahead of the workers, limited by the size of the ring.}
\label{fig:pat-roles}
\end{figure}

\begin{algorithm}[t]
\caption{\texttt{runBineAgPat}, kernel of one channel.}
\label{alg:pat-kernel}
\begin{algorithmic}[1]
\State $(\mathit{count},\mathit{channelOffset},\mathit{channelCount},\mathit{chunkCount})\gets$ \Call{CollCbdPart}{}
\State clear the ring \texttt{patSteps}; \textbf{barrier}
\If{thread is the schedule thread}
  \State $\mathit{geom}\gets$ \Call{Init}{$s,\mathit{count},\mathit{channelOffset},\mathit{channelCount},\mathit{chunkCount},E_s$}; $\mathit{algo}\gets$ \Call{MakeAlgo}{$\mathit{geom}$}
  \Repeat
    \State wait until the entry of the ring is free
    \State $\mathit{algo}.\texttt{getNextOp}$(entry) \Comment{publishes the entry}
  \Until{the entry is the last}
\ElsIf{thread is a worker}
  \State build \texttt{Primitives} in mode \texttt{primsModeBineAg}
  \Repeat
    \State wait until the entry of the ring is published
    \State \texttt{patCopy}(entry); release the entry
  \Until{the entry is the last}
\EndIf
\State \Call{PostFn}{} \Comment{executed by all the threads}
\end{algorithmic}
\end{algorithm}

\paragraph{Differences from \pat{}.}
The ring of descriptors, the number of workers, the polling loops and the call to \texttt{patCopy} are those of the native \pat{} allgather (\cref{sec:nccl-pat-impl}).
Three things change.
\begin{itemize}
  \item The generator: \texttt{PatAGAlgorithm} is replaced by the generator of the strategy.
  \item The mode of the primitives: \texttt{primsModePatAg} becomes \texttt{primsModeBineAg} (\cref{sec:pat-prims}).
  \item The structure of the function: the body is a template parametrised by the generator and by the post hook, shared by the four strategies.
\end{itemize}
A single group of workers is used, since the generators return a parallel factor of $1$: all the operations of a channel are executed in order by the same threads, as the order of the connections requires.
With \texttt{ProtoSimple<1,1>} a chunk is one slice and one step of the connection.

\paragraph{Dispatch.}
A thin wrapper per strategy (\texttt{runBineSendPat}, \texttt{runBinePermutationPat}, \texttt{runBineBlockByBlockPat}, \texttt{runBineDoubleSendPat}) reads \texttt{channel.bine} (\texttt{nDoublingSteps}, the tables) and gives \texttt{runBineAgPat} the lambda that builds the generator.
The strategy is selected at run time with the field \texttt{bufferManagement}, as in \cref{sec:bfly-collectives}.

%==============================================================================
\section{Changes to the \texttt{Simple} primitives}\label{sec:pat-prims}

All the device-side changes are in the \texttt{Simple} protocol (\texttt{prims\_simple.h}).

\paragraph{A new mode.}
The enumeration \texttt{primsMode} gets the value \texttt{primsModeBineAg}, next to the two of \pat{}.
The constructor treats it as a \pat{} mode (flag \texttt{PatMode}: five roles, one thread each, state in shared memory, \texttt{patBarrier}) and differs only in how the peers are computed.
In \pat{}, thread $d$ loads the peer at distance $2^d$.
Here thread $d<s$ loads the connector of dimension $d$ with the partner of the step: from $\Pi$ through \texttt{bineSendPartner} for send, permute and block-by-block, from $\Pi^{\mathrm{dh}}$ (\texttt{dhlvPartners}) for double send.
Receive peer and send peer are the same rank (\cref{lst:pat-peers}).
\begin{lstlisting}[language=C++,caption={Peers of the \texttt{primsModeBineAg} mode (simplified).},label={lst:pat-peers}]
if (tid < S)
  recvPeer = sendPeer = isDoubleSend
      ? bine->dhlvPartners[rank * S + tid]
      : bineSendPartner(bine->partners, rank, tid, S);
else if (hasPreStep && tid == S) {      // strategy "send" only
  BineAgPre pre = bineAgPre(partners, index, order, rank, S);
  if (pre.recvDim == S) recvPeer = pre.recvPeer;
  if (pre.sendDim == S) sendPeer = pre.sendPeer;
}
\end{lstlisting}

\paragraph{Pre-step peers.}
The pre-step of \emph{send} exchanges data with two ranks that are not, in general, partners of the butterfly: the rank that gets my block, $\nu^{-1}(\mathrm{rev}(r))$, and the one that gives me mine, $\mathrm{pos}(r)$.
\texttt{bineAgPre} returns them with their dimension.
If such a peer coincides with the partner of a step it takes that step's dimension, to avoid two connectors with two step counters for one connection; otherwise it takes the extra dimension $s$, served by one more thread.
A rank thus has at most $s+2$ distinct peers, all set up with the other \bine{} connections (\cref{sec:nccl-channels}).

\paragraph{Two flags of \texttt{patCopy}.}
\texttt{patCopy} handles an operation whose source is the input buffer, with a local copy to the output buffer.
The \bine{} allgather needs two more behaviours, selected by flags of the entry.
\begin{itemize}
  \item \texttt{PatSrcOutput}: the data to send is read from the output buffer. It is the case of every send of the butterfly, whose blocks (the own one and those received before) are already in \texttt{recvbuff}.
  \item \texttt{PatNoLocalCopy}: the data is sent without copying it to the output buffer. It is used by the exchange of the pre-step, which must not write slot $r$ of the own rank, and implicitly by \texttt{PatSrcOutput}.
\end{itemize}
When the source is the output buffer, or no local copy is requested, the redundant destination is dropped from the list of the \texttt{reduceCopy}, and the operation is a plain read into the slot.

\paragraph{Receive.}
A received piece is copied from its offset in the slot (\texttt{recvOffset}) to the output buffer at \texttt{outIx}.
No data is forwarded from the slot: each operation either sends or receives, and what must be sent later is read from the output buffer.
This costs one local copy per piece, the price of the arbitrary position of the pieces, and it is not on the critical path of the network.

\paragraph{Limits.}
The following limits are inherited from \pat{}.
The native allgather exists only for \texttt{Simple}, and the \pat{} modes are branches of its specialization of \texttt{Primitives}, so there is no \texttt{LL} or \texttt{LL128} variant.
The native kernel instantiates \texttt{ProtoSimple<1,1>} with \texttt{FanSymmetric}$\langle1\rangle$ and \texttt{Direct}$=0$, and so does this one, so there is no zero-copy path either.
Supporting other protocols would require porting the \pat{} modes to their primitives, which is outside the scope of this thesis.
The code is compiled only for architectures with \texttt{\_\_CUDA\_ARCH\_\_}$\ge600$, which provide the block-scope atomics used by the ring.

%==============================================================================
\section{Post-step of permute}\label{sec:pat-permute}

The strategy \emph{permute} runs the butterfly in the position space, so afterwards the block of rank $b$ is in slot $\mathrm{pos}(b)$ and must go to slot $b$.
The function \texttt{bineInPlacePermute} does it with no scratch buffer, restricted to the columns of the slice of the channel.
Since the implementation uses slice mode, the columns of different channels are disjoint and no barrier among channels is needed, which is what prevents permute from using pos mode in \cref{ch:butterfly}.

For a permutation $f$ of $[0,p-1]$ (here $f(b)=\mathrm{pos}(b)$), the function replaces $\texttt{new}[x]=\texttt{old}[f(x)]$ following the cycles of $f$.
\begin{itemize}
  \item Each cycle is handled by its smallest index, found by walking the cycle from every $i$ and stopping at a smaller element; fixed points are skipped.
  \item For the leader $i$ of a cycle, the threads split the \texttt{channelCount} columns; each saves the element of block $i$, moves along the cycle copying $\texttt{old}[f(\mathit{cur})]$ into $\mathit{cur}$, and closes the cycle with the saved element.
\end{itemize}
The columns are independent, so no synchronization is needed inside the permutation.
The function runs in the hook \texttt{post}, executed by all the threads after the schedule thread and the workers have finished; the hook begins with a \texttt{\_\_syncthreads} so that every received piece is visible, and does nothing if the channel has no data.
Finding the leaders costs $O(\sum_C \lvert C\rvert^2)$ over the cycles $C$ of $f$ for each block, independent of the data, and negligible for the sizes in which \emph{permute} is used.

%==============================================================================
\section{Experimental results}\label{sec:pat-exp}

\todo{Experimental results: comparison with the kernels of \cref{ch:butterfly} (slice and pos mode), the ring and the native \pat{}, for the four strategies.}
